% 示例：霍尔效应及其应用（深圳大学《大学物理实验（2）》格式）
% 数据为一次真实实验记录；封面个人信息已替换为占位。需求清单见 requirements.md。
% 编译：python analysis.py && python ../../skills/lab-report-tex/scripts/build.py report.tex --thumbs preview
\documentclass{labreport}
\def\FitIsK{1.185}%  fit_plot.py fit_is.csv y~x
\def\FitIsB{-0.007048}%  fit_plot.py fit_is.csv y~x
\def\FitIsR{1.0000}%  fit_plot.py fit_is.csv y~x
\def\FitIsUK{8.2e-04}%  fit_plot.py fit_is.csv y~x
\def\FitIsUB{0.002}%  fit_plot.py fit_is.csv y~x
\def\FitImK{11.82}%  fit_plot.py fit_im.csv y~x
\def\FitImB{4.000e-04}%  fit_plot.py fit_im.csv y~x
\def\FitImR{1.0000}%  fit_plot.py fit_im.csv y~x
\def\FitImUK{0.015}%  fit_plot.py fit_im.csv y~x
\def\FitImUB{0.005}%  fit_plot.py fit_im.csv y~x
\def\KH{177}%  analysis.py
\def\XC{110}%  analysis.py
\def\VHC{1.970}%  analysis.py
\def\DV{0.020}%  analysis.py
\def\UVH{0.006}%  analysis.py
\def\RelVH{0.3}%  analysis.py
\def\RelKH{1.2}%  analysis.py
\def\RelIs{0.5}%  analysis.py
\def\RelB{1.3}%  analysis.py
\def\BVal{3.71}%  analysis.py
\def\BU{0.05}%  analysis.py
\def\BLeft{0.438}%  analysis.py
\def\BRight{0.852}%  analysis.py
\def\XLeft{0}%  analysis.py
\def\XRight{200}%  analysis.py
\def\EndRatio{1.9}%  analysis.py
\def\ExIs{3.50}%  analysis.py
\def\ExVa{3.74}%  analysis.py
\def\ExVb{-4.54}%  analysis.py
\def\ExVc{4.54}%  analysis.py
\def\ExVd{-3.75}%  analysis.py
\def\ExVH{4.143}%  analysis.py
\def\ExSum{16.57}%  analysis.py


\labset{
  school     = {深圳大学},
  cover      = {szu},
  coursename = {大学物理实验（2）},
  coursecode = {1000001-00},
  expname    = {霍尔效应及其应用},
  college    = {物理与电子科学学院},
  teacher    = {李老师},
  author     = {张三},
  studentid  = {2025000000},
  group      = {1},
  location   = {实验楼A302},
  date       = {2026年9月23日},
  submitdate = {2026年9月30日},
}

\begin{document}
\makecover

\begin{labsection}{实验目的}
1．了解产生霍尔效应的物理原理；\par
2．学习控制变量法和对称测量法；\par
3．测量霍尔器件的输出特性；\par
4．利用霍尔元件测量长直螺线管的轴向磁场分布。
\end{labsection}

\begin{labsection}{实验原理}
\labsub{霍尔效应}
把载流半导体薄片置于与电流方向垂直的磁场中，载流子受洛伦兹力作用发生横向偏转，在薄片两侧积累电荷并建立横向电场。洛伦兹力与横向电场力平衡时，两侧形成稳定的霍尔电压
\begin{equation}
  V_H = K_H I_s B \label{eq:vh}
\end{equation}
式中 $I_s$ 为工作电流（mA），$B$ 为磁感应强度（T），$K_H$ 为霍尔器件灵敏度（mV/(mA$\cdot$T)），由材料与几何尺寸决定。

\labsub{附加电动势与对称测量法}
实测电压中还叠加有不等电位差、爱廷豪森效应、能斯特效应和里纪-勒杜克效应等附加电动势。分别测量 $(+B,+I_s)$、$(-B,+I_s)$、$(-B,-I_s)$、$(+B,-I_s)$ 四种组合下的电压 $V_1\sim V_4$，取
\begin{equation}
  V_H = \frac{V_1 - V_2 + V_3 - V_4}{4} \label{eq:sym}
\end{equation}
可消除大部分附加电动势。

\labsub{用霍尔元件测磁场}
由式 \eqref{eq:vh}，已知 $K_H$，测出 $V_H$ 与 $I_s$ 即可求得
\begin{equation}
  B = \frac{V_H}{K_H I_s} \label{eq:b}
\end{equation}
\end{labsection}

\begin{labsection}{实验仪器}
霍尔效应实验仪（含霍尔器件、工作电流源、励磁电流源、数字毫伏表）、长直螺线管及位移标尺。主要参数如下：
\labtable{主要仪器及参数}[tab:instr]{
\begin{tabularx}{\linewidth}{|>{\centering\arraybackslash}X|c|c|>{\centering\arraybackslash}X|}
\hline 仪器名称 & 量程 & 分度值 & 仪器误差限（按说明书估计） \\ \hline
数字毫伏表（mV 档） & 0～20\,mV & 0.01\,mV & 0.5\%×读数 + 1 个字 \\ \hline
工作电流源 $I_s$ & 0～10\,mA & 0.01\,mA & 0.5\%×读数 + 1 个字 \\ \hline
励磁电流源 $I_m$ & 0～1\,A & 0.001\,A & —— \\ \hline
螺线管位移标尺 & 0～250\,mm & 1\,mm & —— \\ \hline
霍尔器件 & \multicolumn{2}{c|}{$K_H=\KH$\,mV/(mA$\cdot$T)} & 标称误差 ±2\% \\ \hline
\end{tabularx}}
\end{labsection}

\begin{labsection}{实验内容与步骤}
1．\textbf{测绘 $V_H$-$I_s$ 曲线}：先将霍尔电压调零，取 $I_m=0.300$\,A，工作电流 $I_s$ 依次取 1.00、1.50 … 3.50\,mA，每个 $I_s$ 测四种方向组合下的 $V_1\sim V_4$。\par  % lrt:literal（工作电流的测量设定点，不是计算结果）
2．\textbf{测绘 $V_H$-$I_m$ 曲线}：保持 $I_s=3.00$\,mA，$I_m$ 依次取 0.100～0.500\,A（间隔 0.100\,A），同法测量。\par
3．\textbf{测螺线管轴向磁场}：$I_s=3.00$\,mA、$I_m=0.300$\,A，霍尔片沿轴线从 $X=0$ 移到 200\,mm，每隔 10\,mm 测一组 $V_1\sim V_4$。\par
\textbf{注意事项：}换向时先把电流调到零；霍尔片移动时不要碰到螺线管内壁；工作电流不超过 10\,mA。
\end{labsection}

\begin{labsection}{数据记录和处理}
\textbf{1．霍尔器件输出特性}\par
按式 \eqref{eq:sym} 计算各组 $V_H$（表 \ref{tab:is}、表 \ref{tab:im}）。示例：$I_s=\ExIs$\,mA 时，
$V_H=[\ExVa-(\ExVb)+\ExVc-(\ExVd)]/4=\ExSum/4=\ExVH$\,mV。
% 由 csv2tex.py 从 t_is.csv 生成，改数据请改 CSV 后重新生成
\labtable{$V_H$-$I_s$ 关系测量数据（$I_m=0.300$\,A）}[tab:is]{
\begin{tabular}{cccccc}
\toprule
$I_s$/mA & $V_1$/mV & $V_2$/mV & $V_3$/mV & $V_4$/mV & $V_H$/mV \\
\midrule
1.00 & 1.06 & $-1.30$ & 1.29 & $-1.07$ & 1.180 \\
1.50 & 1.60 & $-1.94$ & 1.94 & $-1.60$ & 1.770 \\
2.00 & 2.13 & $-2.59$ & 2.59 & $-2.14$ & 2.362 \\
2.50 & 2.67 & $-3.24$ & 3.24 & $-2.68$ & 2.958 \\
3.00 & 3.20 & $-3.89$ & 3.89 & $-3.21$ & 3.548 \\
3.50 & 3.74 & $-4.54$ & 4.54 & $-3.75$ & 4.143 \\
\bottomrule
\end{tabular}}

% 由 csv2tex.py 从 t_im.csv 生成，改数据请改 CSV 后重新生成
\labtable{$V_H$-$I_m$ 关系测量数据（$I_s=3.00$\,mA）}[tab:im]{
\begin{tabular}{cccccc}
\toprule
$I_m$/A & $V_1$/mV & $V_2$/mV & $V_3$/mV & $V_4$/mV & $V_H$/mV \\
\midrule
0.100 & 0.85 & $-1.53$ & 1.52 & $-0.85$ & 1.188 \\
0.200 & 2.02 & $-2.70$ & 2.70 & $-2.02$ & 2.360 \\
0.300 & 3.20 & $-3.89$ & 3.89 & $-3.20$ & 3.545 \\
0.400 & 4.38 & $-5.08$ & 5.08 & $-4.38$ & 4.730 \\
0.500 & 5.56 & $-6.27$ & 6.27 & $-5.56$ & 5.915 \\
\bottomrule
\end{tabular}}

\labfigs{fig/vh_is.pdf}{$V_H$-$I_s$ 关系曲线}{fig:is}{fig/vh_im.pdf}{$V_H$-$I_m$ 关系曲线}{fig:im}
最小二乘拟合：$V_H$-$I_s$ 斜率 $k_1=\num{\FitIsK}$\,mV/mA，截距 $\num{\FitIsB}$\,mV，$r=\FitIsR$；
$V_H$-$I_m$ 斜率 $k_2=\num{\FitImK}$\,mV/A，截距 $\num{\FitImB}$\,mV，$r=\FitImR$。

\textbf{2．螺线管轴向磁场分布}\par
由式 \eqref{eq:b} 计算各点磁感应强度（表 \ref{tab:x}，图 \ref{fig:bx}）。示例：$X=\XC$\,mm 处 $V_H=\VHC$\,mV，
$B=\VHC/(\KH\times3.00)\,\mathrm{T}=\BVal$\,mT。
% 由 csv2tex.py 从 t_x.csv 生成，改数据请改 CSV 后重新生成
\labtable{螺线管轴向磁场分布（$I_s=3.00$\,mA，$I_m=0.300$\,A，$K_H=177$\,mV/(mA$\cdot$T)）}[tab:x]{
\begin{tabular}{ccccccc}
\toprule
$X$/mm & $V_1$/mV & $V_2$/mV & $V_3$/mV & $V_4$/mV & $V_H$/mV & $B$/mT \\
\midrule
0 & 0.16 & $-0.31$ & 0.31 & $-0.15$ & 0.233 & 0.438 \\
10 & 0.37 & $-0.48$ & 0.48 & $-0.36$ & 0.422 & 0.796 \\
20 & 0.68 & $-0.81$ & 0.81 & $-0.68$ & 0.745 & 1.403 \\
30 & 1.13 & $-1.22$ & 1.23 & $-1.12$ & 1.175 & 2.213 \\
40 & 1.47 & $-1.57$ & 1.58 & $-1.46$ & 1.520 & 2.863 \\
50 & 1.65 & $-1.77$ & 1.76 & $-1.64$ & 1.705 & 3.211 \\
60 & 1.75 & $-1.85$ & 1.85 & $-1.75$ & 1.800 & 3.390 \\
70 & 1.82 & $-1.92$ & 1.92 & $-1.82$ & 1.870 & 3.522 \\
80 & 1.87 & $-1.97$ & 1.97 & $-1.87$ & 1.920 & 3.616 \\
90 & 1.90 & $-2.00$ & 2.00 & $-1.90$ & 1.950 & 3.672 \\
100 & 1.92 & $-2.02$ & 2.02 & $-1.91$ & 1.968 & 3.705 \\
110 & 1.92 & $-2.02$ & 2.03 & $-1.91$ & 1.970 & 3.710 \\
120 & 1.90 & $-2.00$ & 2.00 & $-1.90$ & 1.950 & 3.672 \\
130 & 1.87 & $-1.97$ & 1.98 & $-1.87$ & 1.923 & 3.621 \\
140 & 1.82 & $-1.92$ & 1.91 & $-1.81$ & 1.865 & 3.512 \\
150 & 1.75 & $-1.85$ & 1.84 & $-1.75$ & 1.798 & 3.385 \\
160 & 1.64 & $-1.74$ & 1.74 & $-1.64$ & 1.690 & 3.183 \\
170 & 1.47 & $-1.57$ & 1.58 & $-1.47$ & 1.522 & 2.867 \\
180 & 1.20 & $-1.30$ & 1.30 & $-1.20$ & 1.250 & 2.354 \\
190 & 0.82 & $-0.92$ & 0.92 & $-0.82$ & 0.870 & 1.638 \\
200 & 0.40 & $-0.50$ & 0.50 & $-0.41$ & 0.452 & 0.852 \\
\bottomrule
\end{tabular}}

\labfig[0.72\linewidth]{fig/b_x.pdf}{螺线管轴向磁场 $B$-$X$ 分布曲线}[fig:bx]

\textbf{3．中心点磁场的不确定度}\par
$V_1\sim V_4$ 为单次测量，只计 B 类。对 $V_H\approx\VHC$\,mV，各电压分量 $\Delta_\text{仪}\approx\DV$\,mV，$u_B=\Delta_\text{仪}/\sqrt3$。由式 \eqref{eq:sym}，
\[
  u(V_H)=\frac{\sqrt{4u_B^2}}{4}=\frac{u_B}{2}\approx\UVH\ \mathrm{mV},\qquad \frac{u(V_H)}{V_H}=\RelVH\%
\]
$K_H$ 标称误差 ±2\%，$u(K_H)/K_H=2\%/\sqrt3=\RelKH\%$；$I_s=3.00$\,mA 时 $u(I_s)/I_s=\RelIs\%$。式 \eqref{eq:b} 为积商形式：
\[
  \frac{u(B)}{B}=\sqrt{\left(\frac{u(V_H)}{V_H}\right)^2+\left(\frac{u(K_H)}{K_H}\right)^2+\left(\frac{u(I_s)}{I_s}\right)^2}=\RelB\%
\]
中心点磁场：$B=(\BVal\pm\BU)$\,mT（$P=0.683$）。
\end{labsection}

\begin{labsection}{结果陈述及总结}
1．由表 \ref{tab:is} 及图 \ref{fig:is}，$V_H$ 随 $I_s$ 线性增大，斜率 $\num{\FitIsK}$\,mV/mA，$r=\FitIsR$，截距接近零，即 $V_H\propto I_s$，与式 \eqref{eq:vh} 一致。\par
2．由表 \ref{tab:im} 及图 \ref{fig:im}，$V_H$ 随 $I_m$ 线性增大，斜率 $\num{\FitImK}$\,mV/A，$r=\FitImR$。螺线管内 $B\propto I_m$，故 $V_H\propto B$，霍尔效应测磁场的原理得到验证。\par
3．由表 \ref{tab:x} 及图 \ref{fig:bx}，轴线上磁场中间强、两端弱：$X\approx100\sim110$\,mm 处最大，$B=(\BVal\pm\BU)$\,mT；$X=\XLeft$\,mm 端 $B=\BLeft$\,mT，$X=\XRight$\,mm 端 $B=\BRight$\,mT，两端相差约 \EndRatio\ 倍，分布并不严格对称。可能原因是位移标尺零点与螺线管端口没有对齐，测量区间相对螺线管中心偏向一侧。\par
4．误差主要来自 $K_H$ 的标称误差（贡献 \RelKH\% ，在 \RelB\% 的合成不确定度中占主导），其次是工作电流读数；对称测量法已消除大部分附加电动势。
\end{labsection}

\begin{labsection}{思考题}
\question{霍尔电压的大小与哪些因素有关？}
由 $V_H=K_HI_sB$，霍尔电压与工作电流 $I_s$、磁感应强度 $B$ 成正比；灵敏度 $K_H$ 由材料（载流子类型与浓度）和厚度决定，并随温度变化。本实验中增大 $I_s$ 或 $I_m$ 时 $V_H$ 均线性增大（$r=\FitIsR$、$\FitImR$），与理论一致。

\question{若已知霍尔器件的性能参数，用霍尔效应法测量一个未知磁场时，测量误差有哪些来源？}
① $V_H$ 的测量误差（毫伏表仪器误差与读数误差）；② 对称测量法未完全消除的附加电动势；③ 工作电流 $I_s$ 的误差及波动；④ $K_H$ 的标称误差与温漂——本实验中它是最大的一项（\RelKH\%）；⑤ 霍尔片平面与磁场不垂直，实际测到 $B\cos\theta$；⑥ 霍尔片的定位误差，在磁场变化快的螺线管端部影响尤其大；⑦ 磁场本身的不均匀性。
\end{labsection}

\labcomment
\labgrading{预习/20,操作及记录/40,数据处理与结果陈述/30,思考题/10}

% 示例没有附真实手写记录：这是一张标明“演示占位图，非真实实验记录”的图片
\labphotos{原始数据（演示占位图，非真实记录）}{raw/demo_placeholder.jpg}
\end{document}
